\chapter{The Seam}

The loop has closed, and a closed loop is a thing you can go around more than once. The machine, having joined its
end to its beginning, is now built to run the whole circuit again --- and again --- and the next question is how
one pass around the loop hands off to the next. There has to be a seam: a place where the pass just
finished gives way to the pass about to begin, where what the last circuit left behind becomes what the next
circuit starts from. The seam is where the machine's loop becomes not a single closed circle but a spiral it can
trace as many times as it is driven to, and the count of those tracings is going to turn out to matter more than
anything the machine has counted yet.

\section{One Pass Seeds the Next}

Each pass around the loop is constructed from the residue the previous pass left, so the circuit can be run over
and over, each turn seeded by the one before it.

Against the machine, the seam is another stretch of construction, inhabited rather than proved, like the descent
of the backward walk. At the join, every rung of the tower is re-instantiated --- built again, fresh, for the new
pass --- but not from nothing: from the residue the completed pass deposited. Recall the leftover the machine
learned to carry back in the finite approach, the gap between an approximation and the number it reached for, held
deliberately rather than discarded. Here that habit pays its structural due. What one pass around the loop does
not finish, it leaves as a residue at the seam, and the next pass picks that residue up as its own starting
condition and carries the circuit forward from there. Pass one closes and deposits what it could not complete;
pass two opens on exactly that deposit and takes the loop around again; and the seam between them is the piece of
construction that makes the hand-off clean --- capabilities reset, the residue forwarded, the machine ready to run
the circuit once more. None of this is a theorem. It is built: the seam is inhabited, rung by rung, as the join
that lets one turn become the seed of the next.

\section{Charge Is the Loop Count}

Now the reading names something carried to the very end. The machine, running its loop again and again, keeps a
count: how many times it has been around. Each pass is one turn, and the turns accumulate, and the accumulated
count of turns is a definite quantity the construction holds. That quantity gets a name, and the name is
\emph{charge}. Charge, in this instrument, is the count of the loop's turns --- the number of
times the circuit has been driven around, tallied at the seam, one more with every pass.

Be exact about the grade, because the word carries weight the construction must be careful to earn. What is
\emph{constructed} is the seam and the counting --- the machine really does re-instantiate each pass from the
last and really does keep a running tally of the turns, and that is inhabited, built, no theorem needed. That the
tally deserves the name \emph{charge} is a reading laid over the built count --- an interpretation, offered and to
be judged, exactly the kind of thing graded as a reading rather than a proof. But it is a reading with
firm ground under it. The loop count is a genuine construction fact: the turns are really counted, the tally is
really kept, and to call that tally the charge is not to smuggle in a physics the machine has not built but to
recognize, in the count of turns, the quantity the later descent reads off records as charge. The count is
the world's, in the seam's own terms --- the machine tallies the turns it is actually driven around, not turns it
invents --- and the name is the hand's, hung on the count the way every reading hangs a name on what the
instrument has built. Charge is the loop count: constructed as a tally, read as a charge, and the two kept
honestly apart.

There is a plain intuition in it worth keeping, though the machine does not cash it here: what a thing owes, in this
machine, is a matter of how many times it has been taken around --- the count of turns is the account, and the
account is what the later physics will bill against. The seam is where that account is kept. Every pass adds one
to it; nothing subtracts; the loop count only rises as the circuit is driven, a tally forwarded across every seam
and never reset. What the machine will eventually make of that rising count --- how it reads it as a physical
charge on a physical record --- waits until the physics can be read off the count. Here the machine only
builds the seam, keeps the tally, and names it: charge is how many times around.
